\section{Introduction}
\label{sec:introduction}

Federated learning (FL) allows institutions or mobile devices to train a shared model collaboratively without ever uploading raw private data~\cite{mcmahan2017communication,kairouz2019advances}. In practice, however, this privacy advantage leaves the central server flying partly blind. Client datasets are often strongly non-IID---one hospital or user may hold rare classes that others rarely see---while a fraction of participants may be compromised and send poisoned updates~\cite{li2020fedprox,kairouz2019advances}. Because the server typically possesses no clean validation dataset of its own~\cite{cao2020fltrust}, it must decide whom to trust using only the updates and reports arriving from the clients themselves.

Most existing defenses tackle this challenge with static geometric filters such as Krum, Trimmed-Mean, or geometric-median aggregation~\cite{blanchard2017machine,yin2018byzantine,pillutla2022robust}. Intuitively, these rules treat updates far from the crowd as suspicious. That intuition works against crude attacks, but creates a fundamental \textbf{robustness--heterogeneity dilemma} under realistic non-IID data. When a filter is tightened enough to stop stealthy, boundary-hugging attacks such as ALIE or Min-Max~\cite{shejwalkar2021manipulating}---which carefully craft poisoned updates to sit just inside the benign cloud---it begins misclassifying \emph{honest minority clients} as outliers. Over dozens of rounds, their aggregation weights shrink toward zero, starving rare classes of representation in the global model. Loosening the filter rescues those minority clients, yet immediately opens the door to adaptive poisoning.

A natural response is to replace static thresholds with an adaptive controller, such as a reinforcement learning (RL) policy that adjusts filtering strictness online. Existing RL-based defenses, however, rely on clean server-side proxy datasets to score each round~\cite{he2025fedaa}, which violates the data-free premise of FL. Conversely, utility-aware schemes that ask clients to report their own training loss reductions~\cite{lai2021oort,fu2022clientselection} can be fooled by \emph{Fake Utility} attacks~\cite{gao2021fifl,xiong2023truthful,chakarov2024incentivizing}, where adversaries inflate their reported usefulness to hijack larger aggregation weights.

Without server-side clean data, closing the loop requires confronting a two-stage threat: adversaries can corrupt their \emph{forward} model updates and utility claims, and can also lie in any \emph{backward} validation feedback used to guide the controller. At first glance, filtering out evaluators whose reported loss improvements differ from the cohort median seems like an obvious fix. Yet in heterogeneous FL, median-residual filtering fails precisely when we need it most. Consider a round where an overly strict filter helps the $60\%$ honest majority while hurting a $20\%$ honest minority; if $20\%$ Byzantine evaluators echo the majority's positive score, the honest minority's negative alerts sit alone in the tail and are discarded as ``anomalies,'' leaving the controller blind to minority harm.

To break this deadlock, we propose \textbf{FedTC} (Federated Trust Calibration), a server-data-free framework that couples an analytically grounded \emph{Forward Calibration Layer} with a minority-preserving \emph{Backward Confidence-Interval Controller}. On the forward pass, FedTC replaces brittle select-or-drop rules with continuous credibility scores anchored at the geometric median and normalized by the cohort's typical spread. Under explicit parameter bounds, we prove that central honest clients never collapse to zero weight, that inflated utility claims cannot breach an exact analytical cutoff radius $R_t(\tau_t)$, and that a clean separating threshold exists whenever honest and malicious deviations are separated.

On the backward pass, FedTC updates its calibration parameters once every $H$ rounds by querying a decoupled evaluator cohort for bounded local validation loss changes. Instead of discarding evaluators that disagree with the median, FedTC applies worst-case upper- and lower-tail trimming calibrated to a high-probability malicious budget, yielding a confidence interval $[L_n, U_n]$ that provably contains the true macro-average honest improvement $\bar{Z}_n^\star$. Because negative alerts from damaged minority clients remain intact in $L_n$, over-filtering actions are penalized directly. Coupled with PopArt-normalized dual interval critics and an execution-level fallback gate against a conservative reference policy, FedTC explores permissive, utility-guided calibration only when backed by lower-bound evidence, and safely reverts otherwise. The remainder of the paper formalizes the threat model (Section~\ref{sec:preliminaries}), details FedTC's forward--backward design and guarantees (Section~\ref{sec:method}), and outlines our empirical evaluation and related work (Sections~\ref{sec:experiments}--\ref{sec:related_work}).
